Queuing network control is essential for managing congestion in job-processing systems 
such as service systems, communication networks, and manufacturing processes. Despite growing interest in applying reinforcement learning (RL) techniques,  queueing network control
%model-free reinforcement learning (RL) have been proposed to train control policies for arbitrary network topologies. 
poses distinct challenges, including high stochasticity, large state and action spaces, and lack of stability. 
To tackle these challenges, we propose a scalable framework for policy optimization based on differentiable discrete event simulation. 
Our main insight is that by implementing a well-designed smoothing technique for discrete event dynamics, we can compute $\pathwise$ policy gradients for large-scale queueing networks using auto-differentiation software (e.g., Tensorflow, PyTorch) and GPU parallelization. Through extensive empirical experiments, we observe that our policy gradient estimators are several orders of magnitude more accurate than typical $\reinforce$-based  estimators. 
In addition, we propose a new policy architecture, which drastically improves stability while maintaining the flexibility of neural-network policies.
In a wide variety of scheduling and admission control tasks, we demonstrate that training control policies with pathwise gradients leads to a 50-1000x improvement in sample efficiency over state-of-the-art RL methods.
Unlike prior tailored approaches to queueing, our methods can flexibly handle realistic scenarios, including systems operating in non-stationary environments and those with non-exponential interarrival/service times.


% To tackle these challenges, 
% we propose a scalable framework for policy optimization based on differentiable simulation
% and demonstrate substantial improvements in  sample efficiency in multi-class queuing networks.
% %and reduces sample inefficiency of existing learning approaches. 
% Our main algorithmic insight is to compute pathwise gradients for 
% discrete event dynamical systems with auto-differentiation (e.g. Tensorflow or PyTorch).
% We address the non-differentiability of the discrete event dynamics through carefully designed smoothing. 
% We find that for scheduling and admission control tasks, training control policies with pathwise gradients leads to a 50-1000x improvement in sample efficiency over state-of-the-art RL methods.
% Our methods can handle realistic instances, including systems operating in a non-stationary environment and with non-exponential system inputs (i.e., interarrival and service times), which are relevant for many practical settings.\hntodo{This last bit is confusing because everyone who don't know D and G (meaning almost every reader) will be left perplexed why black-box RL methods don't apply. Need to qualify by saying queueing-specific innovations that stabilize SoTA RL methods etc...}


%%% TeX-master: "abstract"
%%% Local Variables: %%% mode: latex %%% TeX-master: "main" %%% End:
